\documentclass[12pt,a4paper]{article}

\usepackage[margin=1in]{geometry}
\usepackage{amsmath}
\usepackage{booktabs}
\usepackage{enumitem}
\usepackage{siunitx}

\title{Lab 4: Constant Strain Triangle (CST) Element and FEM Assembly}
\author{}
\date{}

\begin{document}

\maketitle

\section*{Introduction}

This laboratory exercise introduces the Constant Strain Triangle (CST)
finite element for two-dimensional, linear elastic, plane-stress analysis.
The three problems below demonstrate the formulation of a single CST
element, verification using a patch test, and assembly of a CST model for a
cantilever beam.

Unless otherwise stated, the material is assumed to be linear, isotropic, and
elastic under plane-stress conditions. The calculations use consistent
N--mm--MPa units. Therefore, the beam dimensions given in metres are written
as millimetres in the numerical model.

\section{Problem 1: Single CST Element}

\subsection*{Problem statement}

Analyse one triangular Constant Strain Triangle element using the finite
element method. Calculate the element area, strain--displacement matrix,
material constitutive matrix, element stiffness matrix, nodal displacements,
strains, and stresses. Calculate the support reactions and compare the
computed response with the expected analytical behaviour.

\subsection*{Given values}

\begin{table}[h!]
\centering
\begin{tabular}{ll}
\toprule
\textbf{Quantity} & \textbf{Value} \\
\midrule
Element type & Constant Strain Triangle (CST) \\
Node 1 & $(0,0)\ \si{\milli\metre}$ \\
Node 2 & $(1,0)\ \si{\milli\metre}$ \\
Node 3 & $(0,1)\ \si{\milli\metre}$ \\
Young's modulus, $E$ & $\SI{200000}{\mega\pascal}$ \\
Poisson's ratio, $\nu$ & $0.30$ \\
Thickness, $t$ & $\SI{1}{\milli\metre}$ \\
\bottomrule
\end{tabular}
\end{table}

The boundary conditions are
\[
u_1=v_1=0,\qquad v_2=0.
\]
The applied load is a vertical force of
\[
F_{y3}=-\SI{100}{\newton}
\]
at Node 3.

\section{Problem 2: CST Patch Test}

\subsection*{Problem statement}

Perform a patch test on a unit square divided into two triangular CST
elements. Prescribe a linear displacement field at the four nodes, calculate
the strain in each element, and compare the numerical strain with the exact
strain obtained by differentiating the prescribed displacement field. The
test is used to verify that the CST formulation reproduces a compatible
linear displacement field exactly.

\subsection*{Given values}

\begin{table}[h!]
\centering
\begin{tabular}{ccl}
\toprule
\textbf{Node} & \textbf{Coordinates} & \textbf{Prescribed displacement} \\
\midrule
1 & $(0,0)\ \si{\milli\metre}$ & $(u_1,v_1)=(0,0)\ \si{\milli\metre}$ \\
2 & $(1,0)\ \si{\milli\metre}$ & $(u_2,v_2)=(0.001,0)\ \si{\milli\metre}$ \\
3 & $(1,1)\ \si{\milli\metre}$ & $(u_3,v_3)=(0.001,0.002)\ \si{\milli\metre}$ \\
4 & $(0,1)\ \si{\milli\metre}$ & $(u_4,v_4)=(0,0.002)\ \si{\milli\metre}$ \\
\bottomrule
\end{tabular}
\end{table}

The element connectivity is
\[
\text{Element 1}=(1,2,3),\qquad
\text{Element 2}=(1,3,4).
\]

The prescribed displacement field is
\[
u(x,y)=0.001x,\qquad v(x,y)=0.002y.
\]
Hence, the exact engineering strain is
\[
\boldsymbol{\varepsilon}_{\mathrm{exact}}
 =
\begin{bmatrix}
\varepsilon_x\\
\varepsilon_y\\
\gamma_{xy}
\end{bmatrix}
 =
\begin{bmatrix}
0.001\\
0.002\\
0
\end{bmatrix}.
\]

The material properties are
\[
E=\SI{200000}{\mega\pascal},\qquad
\nu=0.30,\qquad
t=\SI{1}{\milli\metre},
\]
with plane-stress behaviour.

\section{Problem 3: Cantilever Beam Using CST Elements}

\subsection*{Problem statement}

Model a rectangular cantilever beam using triangular CST elements. Generate a
structured mesh, divide every rectangular mesh cell into two triangles,
assemble the global stiffness matrix, apply the fixed-end boundary
conditions, and calculate the vertical displacement at the loaded tip.
Repeat the analysis for progressively finer meshes and compare the finite
element tip displacement with the Euler--Bernoulli analytical solution.

\subsection*{Given values}

\begin{table}[h!]
\centering
\begin{tabular}{ll}
\toprule
\textbf{Quantity} & \textbf{Value} \\
\midrule
Beam length, $L$ & $\SI{1}{\metre}=\SI{1000}{\milli\metre}$ \\
Beam height, $H$ & $\SI{0.1}{\metre}=\SI{100}{\milli\metre}$ \\
Thickness, $t$ & $\SI{1}{\milli\metre}$ \\
Young's modulus, $E$ & $\SI{210000}{\mega\pascal}$ \\
Poisson's ratio, $\nu$ & $0.30$ \\
Applied tip load, $F_y$ & $-\SI{100}{\newton}$ \\
Left-edge boundary condition & $u=v=0$ \\
\bottomrule
\end{tabular}
\end{table}

The mesh must contain at least 10 segments along the beam length and 1--2
segments through its height. The mesh-refinement study uses the following
configurations:
\[
(N_x,N_y)=(10,2),\ (20,4),\ (40,8),\ (80,16),
\]
where $N_x$ is the number of longitudinal segments and $N_y$ is the number
of height segments.

For comparison, calculate the second moment of area and analytical tip
displacement from
\[
I=\frac{tH^3}{12},
\qquad
\delta_{\mathrm{tip}}
=\frac{|F_y|L^3}{3EI}.
\]

\section*{Expected discussion}

For each problem, report the calculated matrices and numerical results.
Discuss whether the single-element and patch-test results agree with the
analytical values. For the cantilever, report the mesh-dependent tip
displacement and percentage error. Explain why a coarse CST mesh is
typically too stiff in bending and why mesh refinement improves the result.

\end{document}
